% GENERATED FILE. Edit canonical lessons, then run npm run generate:books.
% canonical-source-hash: fnv1a64:aadf2b8f5787db46
% canonical-lessons: ZH-C72-huoche, ZH-W72-huo, ZH-C72-kaiche, ZH-W72-kai, ZH-C72-feiji, ZH-C72-lu, ZH-C72-dao

\chapter{Getting About}
\label{ch:zh-getting-about}
\hlchaptermodality{72}

\begin{chapteropening}
\textbf{By the end of this chapter:} \emph{I can talk about trains, driving, planes, roads and arriving.}

Complete the last lesson of chapter 72: I can talk about trains, driving, planes, roads and arriving.
\end{chapteropening}

\section[\texorpdfstring{\zh{火车} (huǒchē)}{火车 (huǒchē)}]{\zh{火车} (huǒchē) — train}

\label{lesson:ZH-C72-huoche}



\begin{quote}
Before the new one: say the Mandarin for salty, then the Mandarin for to order food.
\end{quote}

\subsection*{You'll want to know: \zh{火车}}
\textbf{\zh{火车}} — \emph{huǒchē} — \textquotedblleft{}train\textquotedblright{}.

\begin{quote}
\textbf{\zh{这儿的火车很长}} \emph{zhèr de huǒchē hěn cháng} — The trains here are long.
\end{quote}

The new character, in pieces: \zh{火} — a dot on each side, and \zh{人} between them. The next lesson puts it in your hand.

\begin{cousinweb}
\textbf{\zh{火车}} (\emph{huǒchē}) is built from \zh{火} 'fire' and \zh{车} 'vehicle', both inherited from Middle Chinese, \emph{xwaX} and \emph{tsyhae} (Guangyun, in Baxter's transcription): a fire carriage, from the days of steam.
\end{cousinweb}

\begin{grammarlens}[title={What you've built}]
One more word for times and partings.
\end{grammarlens}

\subsection*{Your turn}
\begin{itemize}
\raggedright
  \item \emph{Say it:} \emph{huǒchē}
  \item \emph{Say it:} \emph{huǒchē}, once more
  \item \emph{Say it:} \emph{zhèr de huǒchē hěn cháng} — The trains here are long.
  \item \emph{Recall:} say the Mandarin for salty, then the Mandarin for to order food, then say \emph{huǒchē} again
\end{itemize}

\begin{tcolorbox}[breakable,colback=teal!4,colframe=teal!35!black,title={Before you move on}]
What does \zh{火车} mean? (\textquotedblleft{}Train\textquotedblright{}.) Say it once more.
\end{tcolorbox}
\section[\texorpdfstring{\zh{火} (huǒ)}{火 (huǒ)}]{\zh{火} — two sparks around \zh{人}}

\label{lesson:ZH-W72-huo}



\begin{quote}
The word first, then the character inside it.

\begin{itemize}
\raggedright
  \item \emph{Recall:} say \emph{huǒchē} — \textbf{R1}, one lesson back, before it decays
\end{itemize}

\begin{quote}
\zh{人} \zh{大}
\end{quote}

\emph{Read it:} each one aloud, with nothing said first — \emph{rén}, \emph{dà}
\end{quote}

\subsection*{Script — two sparks around \zh{人}}
\hlblockfigure{\detokenize{figures/ZH-W72-huo-filmstrip.pdf}}{How \zh{火} is written, stroke by stroke}

\begin{quote}
\zh{火}
\end{quote}

\textbf{One piece, four strokes.}

\textbf{A short dot} on the left, slanting down to the right.

\textbf{A short left-falling stroke} on the right, falling toward the middle.

\textbf{\zh{人}, two strokes}: a long left-falling stroke from the top down to the lower left, then a long stroke falling to the right from its middle.

2 + 2 = four.

\subsection*{Your turn}
\begin{itemize}
\raggedright
  \item \emph{Write it:} \zh{火}, all four, with the model in front of you
  \item \emph{Cover:} write it again from memory, and count the pieces, not the strokes
\end{itemize}

\begin{tcolorbox}[breakable,colback=teal!4,colframe=teal!35!black,title={Before you move on}]
How many strokes in \zh{火}? (\textbf{Four}.) Which character you know is its last two strokes? (\textbf{\zh{人}}.) What comes first? (\textbf{The dot on the left}.)

Source: \href{https://raw.githubusercontent.com/chanind/hanzi-writer-data/68d10a4b21150cae5e1ebbd223eed289cf32d90c/data/\%E7\%81\%AB.json}{Hanzi Writer Data \zh{火}, pinned stroke paths and medians}.
\end{tcolorbox}
\section[\texorpdfstring{\zh{开车} (kāichē)}{开车 (kāichē)}]{\zh{开车} (kāichē) — to drive}

\label{lesson:ZH-C72-kaiche}



\begin{quote}
Before the new one: say the Mandarin for to order food, then the Mandarin for train.
\end{quote}

\subsection*{You'll want to know: \zh{开车}}
\textbf{\zh{开车}} — \emph{kāichē} — \textquotedblleft{}to drive\textquotedblright{}.

\begin{quote}
\textbf{\zh{明天我们开车去看山}} \emph{míngtiān wǒmen kāichē qù kàn shān} — Tomorrow we're driving out to see the mountains.
\end{quote}

The new character, in pieces: \zh{开} — two bars, with two legs down through them. The next lesson puts it in your hand.

\begin{cousinweb}
\textbf{\zh{开车}} (\emph{kāichē}) is built from \zh{开} 'to open' and \zh{车} 'vehicle', both inherited from Middle Chinese, \emph{khoj} and \emph{tsyhae} (Guangyun, in Baxter's transcription). Driving is not among the senses of the first there; opening up, setting going, is the likely step, but it is not settled.
\end{cousinweb}

\begin{grammarlens}[title={What you've built}]
One more word for times and partings.
\end{grammarlens}

\subsection*{Your turn}
\begin{itemize}
\raggedright
  \item \emph{Say it:} \emph{kāichē}
  \item \emph{Say it:} \emph{kāichē}, once more
  \item \emph{Say it:} \emph{míngtiān wǒmen kāichē qù kàn shān} — Tomorrow we're driving out to see the mountains.
  \item \emph{Recall:} say the Mandarin for to order food, then the Mandarin for train, then say \emph{kāichē} again
\end{itemize}

\begin{tcolorbox}[breakable,colback=teal!4,colframe=teal!35!black,title={Before you move on}]
What does \zh{开车} mean? (\textquotedblleft{}To drive\textquotedblright{}.) Say it once more.
\end{tcolorbox}
\section[\texorpdfstring{\zh{开} (kāi)}{开 (kāi)}]{\zh{开} — two bars, two legs}

\label{lesson:ZH-W72-kai}



\begin{quote}
The word first, then the character inside it.

\begin{itemize}
\raggedright
  \item \emph{Recall:} say \emph{kāichē} — \textbf{R1}, one lesson back, before it decays
\end{itemize}

\begin{quote}
\zh{二} \zh{车}
\end{quote}

\emph{Read it:} each one aloud, with nothing said first — \emph{èr}, \emph{chē}
\end{quote}

\subsection*{Script — two bars, two legs}
\hlblockfigure{\detokenize{figures/ZH-W72-kai-filmstrip.pdf}}{How \zh{开} is written, stroke by stroke}

\begin{quote}
\zh{开}
\end{quote}

\textbf{One piece, four strokes, the bars first.}

\textbf{A short horizontal} on top, then \textbf{a long horizontal} below it.

\textbf{A left-falling stroke} from the upper bar, down through the lower one to the bottom left.

\textbf{A vertical} on the right, straight down through both bars.

\subsection*{Your turn}
\begin{itemize}
\raggedright
  \item \emph{Write it:} \zh{开}, all four, with the model in front of you
  \item \emph{Cover:} write it again from memory, and count the pieces, not the strokes
\end{itemize}

\begin{tcolorbox}[breakable,colback=teal!4,colframe=teal!35!black,title={Before you move on}]
How many strokes in \zh{开}? (\textbf{Four}.) Which come first, the bars or the legs? (\textbf{The bars}.) Which leg is straight? (\textbf{The right one}.)

Source: \href{https://raw.githubusercontent.com/chanind/hanzi-writer-data/68d10a4b21150cae5e1ebbd223eed289cf32d90c/data/\%E5\%BC\%80.json}{Hanzi Writer Data \zh{开}, pinned stroke paths and medians}.
\end{tcolorbox}
\section[fēijī]{fēijī — aeroplane}

\label{lesson:ZH-C72-feiji}



\begin{quote}
Before the new one: say the Mandarin for train, then the Mandarin for to drive.
\end{quote}

\subsection*{You'll want to know: fēijī}
\emph{fēijī} — \textquotedblleft{}aeroplane\textquotedblright{}.

\begin{quote}
\emph{wǒmen zuò fēijī qù Běijīng} — We're flying to Beijing.
\end{quote}

Its characters come later; for now it is a word for the ear.

\begin{cousinweb}
\emph{fēijī} is built from fēi 'to fly' and \zh{机} 'machine', both inherited from Middle Chinese, \emph{pj+j} and \emph{kj+j} (Guangyun, in Baxter's transcription): a flying machine.
\end{cousinweb}

\begin{grammarlens}[title={What you've built}]
One more word for times and partings.
\end{grammarlens}

\subsection*{Your turn}
\begin{itemize}
\raggedright
  \item \emph{Say it:} \emph{fēijī}
  \item \emph{Say it:} \emph{fēijī}, once more
  \item \emph{Say it:} \emph{wǒmen zuò fēijī qù Běijīng} — We're flying to Beijing.
  \item \emph{Recall:} say the Mandarin for train, then the Mandarin for to drive, then say \emph{fēijī} again
\end{itemize}

\begin{tcolorbox}[breakable,colback=teal!4,colframe=teal!35!black,title={Before you move on}]
What does fēijī mean? (\textquotedblleft{}Aeroplane\textquotedblright{}.) Say it once more.
\end{tcolorbox}
\section[lù]{lù — road; way}

\label{lesson:ZH-C72-lu}



\begin{quote}
Before the new one: say the Mandarin for to drive, then the Mandarin for aeroplane.
\end{quote}

\subsection*{You'll want to know: lù}
\emph{lù} — \textquotedblleft{}road; way\textquotedblright{}.

\begin{quote}
\emph{lù hěn yuǎn} — It's a long way.
\end{quote}

Its characters come later; for now it is a word for the ear.

\begin{cousinweb}
\emph{lù} is inherited from Middle Chinese \emph{luH} (Guangyun, in Baxter's transcription), glossed 'road'.
\end{cousinweb}

\begin{grammarlens}[title={What you've built}]
One more word for times and partings.
\end{grammarlens}

\subsection*{Your turn}
\begin{itemize}
\raggedright
  \item \emph{Say it:} \emph{lù}
  \item \emph{Say it:} \emph{lù}, once more
  \item \emph{Say it:} \emph{lù hěn yuǎn} — It's a long way.
  \item \emph{Recall:} say the Mandarin for to drive, then the Mandarin for aeroplane, then say \emph{lù} again
\end{itemize}

\begin{tcolorbox}[breakable,colback=teal!4,colframe=teal!35!black,title={Before you move on}]
What does lù mean? (\textquotedblleft{}Road; way\textquotedblright{}.) Say it once more.
\end{tcolorbox}
\section[dào]{dào — to arrive; to reach}

\label{lesson:ZH-C72-dao}



\begin{quote}
Before the new one: say the Mandarin for aeroplane, then the Mandarin for road.
\end{quote}

\subsection*{You'll want to know: dào}
\emph{dào} — \textquotedblleft{}to arrive; to reach\textquotedblright{}.

\begin{quote}
\emph{huǒchē hěn màn shí diǎn dào} — The train is slow; it gets in at ten.
\end{quote}

Its characters come later; for now it is a word for the ear.

\begin{cousinweb}
\emph{dào} is inherited from Middle Chinese \emph{tawH} (Guangyun, in Baxter's transcription), glossed 'to reach'.
\end{cousinweb}

\begin{grammarlens}[title={What you've built}]
One more word for times and partings.
\end{grammarlens}

\subsection*{Your turn}
\begin{itemize}
\raggedright
  \item \emph{Say it:} \emph{dào}
  \item \emph{Say it:} \emph{dào}, once more
  \item \emph{Say it:} \emph{huǒchē hěn màn shí diǎn dào} — The train is slow; it gets in at ten.
  \item \emph{Recall:} say the Mandarin for aeroplane, then the Mandarin for road, then say \emph{dào} again
\end{itemize}

\begin{tcolorbox}[breakable,colback=teal!4,colframe=teal!35!black,title={Before you move on}]
What does dào mean? (\textquotedblleft{}To arrive; to reach\textquotedblright{}.) Say it once more.
\end{tcolorbox}
